\section{One-dimensional auxiliary-field formalism and renormalization of the off-lightcone Wilson-line operators}
\label{sec:EFT}
\subsection{The one-dimensional auxiliary-field formalism}
The one-dimensional auxiliary-field formalism is defined by enlarging the QCD Lagrangian to include an extra term \cite{Dorn:1986dt},
\begin{eqnarray}
\label{eq:extendLBare}
\mathcal{L}_{h_v} =\bar{h}_{v, 0} (iv\cdot D) h_{v, 0}, 
\end{eqnarray}
where the subscript ``$0$”  indicates a bare quantity, $h_{v, 0}$ is an auxiliary ``heavy”  Grassmann (or complex) scalar field in either fundamental or adjoint representations of the SU($3$) gauge group, $D_0^{\mu} = \partial^{\mu} - ig_{0}A_0^{\mu, a}T^a$ is the  covariant derivative and $T^a$ are the SU($3$) generators in the appropriate representation. Note that, for $v=(1,0,0,0)$ and $h_{v, 0}$ being in the fundamental representation, the effective Lagrangian in eq.~(\ref{eq:extendLBare}) essentially coincides with the Lagrangian of the heavy-quark effective theory (HQET) in the rest frame \cite{Eichten:1989zv}, except that we choose $h_{v, 0}$ to be a  scalar.

The effective Lagrangian $\mathcal{L}_{h_v}$ is amenable to renormalization through the redefinitions,
\begin{eqnarray}
h_{v, 0} = Z_{h_v}^{\frac{1}{2}} h_v,\, \, \, \, g_0 =Z_g g,\, \, \, \,   A_0 = Z_A^{\frac{1}{2}} A,
\end{eqnarray}
where  $Z_{h_v}$ is the wave-function renormalization constant of the ``heavy” auxiliary field, $Z_g$ and $Z_A$ are the usual QCD renormalization constants of the strong coupling and the gluon field, respectively. There will be a ``mass correction” $i\delta m$ for the ``heavy” auxiliary field and the effective Lagrangian $\mathcal{L}_{h_v} $ expressed in terms of  the renormalized fields and parameters is given by \footnote{The ``mass correction” $i\delta m$ is imaginary for space-like $v$, but real for time-like $v$ in Minkowski spacetime. In Euclidean spacetime,  $i\delta m$ is alway  imaginary. },
\begin{eqnarray}
\label{eq:extendLRen}
\mathcal{L}_{h_v} =Z_{h_v}\bar{h}_v (iv\cdot \partial - i\delta m) h_v + gZ_gZ_A^{\frac{1}{2}}Z_{h_v} \bar{h}_v v\cdot A^{a}T^a h_v,
\end{eqnarray}
where proper representations are implied for the fields and the generator $T^a$. 

The ``mass correction” $i\delta m$ is linearly divergent. The linear divergence vanishes in dimensional regularization, but persists, in lattice regularization and the gradient-flow formalism. Furthermore, the ``mass correction” $i\delta m$ could incorporate contributions of $\mathcal{O}(\Lambda_{\rm QCD})$, attributed to the renormalon ambiguities that are recognized to be present in HQET \cite{Bigi:1994em, Beneke:1994sw}. 

Based on the one-dimensional auxiliary-field formalism, up to a constant normalization factor, the connected two-point function of the ``heavy” auxiliary fields can be related to the Wilson lines through \cite{Dorn:1986dt,Braun:2020ymy,Ji:2020ect},
\begin{eqnarray}
\langle h_{v,0}(x)\bar{h}_{v,0}(0) \rangle_{h_v} = W\left(\frac{v\cdot x}{v^2}, 0\right)\theta\left(\frac{v\cdot x}{v^2}\right)\delta^{(d-1)}(x_\perp),
\end{eqnarray}
where $\langle ... \rangle_{h_v}$ stands for integrating out the ``heavy” auxiliary field $h_v$,  and $x_{\perp}^{\mu} = x^{\mu} - v^{\mu}(v\cdot x)/v^2$.
In this way,  the Wilson-line operators can be substituted with products of local current operators. For the Wilson-line operators defined in eq.~(\ref{eq:defqoperator}) and eq.~(\ref{eq:defgoperator}), we have
\begin{eqnarray}
\mathcal{O}^{\rm B}_{\Gamma}(zv)  &=& \int d^d x\, \delta\left(\frac{v\cdot x}{v^2} -z\right)\langle \bar{\psi}_{0}(x)h_{v,0}(x) \Gamma \bar{h}_{v,0}(0)\psi_0(0)\rangle_{h_v},\\
\mathcal{O}^{\mu\nu\alpha\beta, \rm B}(zv)  &=& \int d^d x\, \delta\left(\frac{v\cdot x}{v^2} -z\right)\langle g^2F_{0}^{\mu\nu}(x)h_{v,0}(x)  \bar{h}_{v,0}(0)F_{0}^{\alpha\beta}(0)\rangle_{h_v},
\end{eqnarray}
where the superscript ``B” indicates bare composite operators, the bare ``heavy” auxiliary fields $h_{v,0}$ in $\mathcal{O}^{\rm B}_{\Gamma}(zv) $, $\mathcal{O}^{\mu\nu\alpha\beta, \rm B}(zv)$ are in fundamental representation and adjoint representation, respectively. 


\subsection{Renormalization of the off-lightcone Wilson-line operators}
The operator as defined in eq.~(\ref{eq:defgoperator}) is not renormalized multiplicatively because, in the presence of the external four-vector $v^{\mu}$,  the components parallel and transverse to $v^{\mu}$ are renormalized differently \cite{Braun:2020ymy}. We introduce the following projectors,
\begin{eqnarray}
g^{\mu\nu}_{\|} = \frac{v_{\mu}v_{\nu}}{v^2},\, \, \, \, g^{\mu\nu}_{\perp} = g^{\mu\nu} - \frac{v_{\mu}v_{\nu}}{v^2},
\end{eqnarray}
to project out the parallel and the transverse components of the field strength tensor $F^{\mu\nu}$ by defining,
\begin{eqnarray}
F^{\mu\nu}_{\|\perp} &=& g^{\mu\alpha}_{\|}g^{\nu\beta}_{\perp}F^{\alpha\beta},\\
F^{\mu\nu}_{\perp\perp} &=& g^{\mu\alpha}_{\perp}g^{\nu\beta}_{\perp}F^{\alpha\beta}.
\end{eqnarray}
Notice that we ignore the mixing with gauge-noninvariant operators since they do not contribute to gauge-invariant observables. Nevertheless, we will verify the mixing patterns found in refs.~\cite{Dorn:1981wa,Zhang:2018diq,Wang:2019tgg} whenever they are relevant for our computations.
Obviously, when $v=(1,0,0,0)$ and $\mu\neq\nu$, $F^{\mu\nu}_{\|\perp}, F^{\mu\nu}_{\perp\perp}$ are proportional to the chromoelectric field \textbf{E} and the chromomagnetic field \textbf{B}, respectively. We introduce the following gluonic Wilson-line operators,
\begin{eqnarray}
\label{eq:defpperp}
\mathcal{O}_{\|\perp}^{\mu\nu\alpha\beta} (zv)&=& g^2 F_{\|\perp}^{\mu\nu}(zv)W(zv,0)F_{\|\perp}^{\alpha\beta}(0),\\
\label{eq:defperpperp}
\mathcal{O}_{\perp\perp}^{\mu\nu\alpha\beta} (zv)&=& g^2 F_{\perp\perp}^{\mu\nu}(zv)W(zv,0)F_{\perp\perp}^{\alpha\beta}(0).
\end{eqnarray}
The off-lightcone Wilson-line operators defined in eq.~(\ref{eq:defqoperator}), eq.~(\ref{eq:defpperp}), and eq.~(\ref{eq:defperpperp})  undergo multiplicative renormalization within the auxiliary-field formalism, as proved in refs.~\cite{Green:2017xeu, Ji:2017oey, Zhang:2018diq}, and shown diagrammatically in refs.~\cite{Ishikawa:2017faj, Li:2018tpe}. That is, the renormalized off-lightcone Wilson-line operators in eq.~(\ref{eq:defqoperator}), eq.~(\ref{eq:defpperp}) and eq.~(\ref{eq:defperpperp}) are related to the bare operators by the formulas,
\begin{eqnarray}
\mathcal{O}^{\rm B}_{\Gamma}(zv, \Lambda) &=& Z_{q}(\Lambda,\mu) e^{\delta m(\Lambda)z} \mathcal{O}^{\rm R}_{\Gamma}(zv, \mu),\\
\mathcal{O}^{\mu\nu\alpha\beta,\, \rm B}_{\|\perp}(zv, \Lambda) &=& Z_{\|\perp}(\Lambda,\mu) e^{\delta m(\Lambda)z} \mathcal{O}_{\|\perp}^{\mu\nu\alpha\beta,\, \rm R}(zv, \mu),\\
\mathcal{O}^{\mu\nu\alpha\beta,\, \rm B}_{\perp\perp}(zv, \Lambda) &=& Z_{\perp\perp}(\Lambda,\mu) e^{\delta m(\Lambda)z} \mathcal{O}_{\|\perp}^{\mu\nu\alpha\beta,\, \rm R}(zv, \mu),
\end{eqnarray}
where the superscript ``R” stands for renormalized composite operators, $\mu$ is the renormalization scale, $\Lambda$ represents a UV cutoff (UV regulator), $\delta m(\Lambda)$ encodes the ``mass correction” of the Wilson line, with its  its  linear divergence, while $Z_{q}(\Lambda,\mu)$, $Z_{\|\perp}(\Lambda,\mu)$ and $Z_{\perp\perp}(\Lambda,\mu)$ encode all logarithmic divergences originating from wave-function and vertex renormalizations.


Within the one-dimensional auxiliary-field formalism, after subtracting linear divergences, the renormalization of the off-lightcone Wilson-line operators defined in eq.~(\ref{eq:defqoperator}), eq.~(\ref{eq:defpperp}), and eq.~(\ref{eq:defperpperp}) is simplified into the renormalization of two local ``heavy-to-light” or ``heavy-to-gluon” currents \footnote{See refs.~\cite{Stefanis:1983ke,Stefanis:2012if} for the renormalization of off-lightcone Wilson-line operators in a different approach based on the Mandelstam formalism.},
\begin{eqnarray}
\label{eq:currentoperator}
J_q(x) &=& \bar{\psi}(x) h_v(x),\\
J^{\mu\nu}_{\|\perp}(x) &=& gF_{\|\perp}^{\mu\nu}(x)h_v(x),\\
J^{\mu\nu}_{\perp\perp}(x) &=& gF_{\perp\perp}^{\mu\nu}(x)h_v(x).
\end{eqnarray}
The renormalization formulas for the above  gauge-invariant local current operators are given by,
\begin{eqnarray}
J_q^{\rm B} (x, \Lambda) &=& Z_{J_q} (\Lambda,\mu)J_q^{\rm R}(x, \mu),\\
J^{\mu\nu,\, \rm B}_{\|\perp}(x, \Lambda) &=& Z_{J,\, \|\perp}(\Lambda,\mu)J^{\mu\nu,\, \rm R}_{\|\perp}(x, \mu),\\
J^{\mu\nu,\, \rm B}_{\perp\perp}(x, \Lambda) &=& Z_{J,\, \perp\perp}(\Lambda,\mu)J^{\mu\nu,\, \rm R}_{\perp\perp}(x, \mu),
\end{eqnarray}
leading to the equivalence 
\begin{eqnarray}
Z_{q} (\Lambda,\mu)&=& Z_{J_q}^2(\Lambda,\mu),\\
Z_{\|\perp} (\Lambda,\mu)&=& Z^2_{J,\, \|\perp}(\Lambda,\mu),\\
Z_{\perp\perp} (\Lambda,\mu)&=& Z^2_{J,\, \perp\perp}(\Lambda,\mu).
\end{eqnarray}
In the following text, we may omit indicating the arguments $\Lambda$ and $\mu$ in the renormalization constants.

The renormalization constants for the local current operators can be further decomposed as
\begin{eqnarray}
Z_{J_q} &=& Z_{\psi}^{\frac{1}{2}} Z_{h_v}^{\frac{1}{2}}Z^V_{J},\\
Z_{J,\, \|\perp} &=& Z_{A}^{\frac{1}{2}} Z_{h_v}^{\frac{1}{2}}Z^V_{\|\perp},\\
Z_{J,\, \perp\perp} &= & Z_{A}^{\frac{1}{2}} Z_{h_v}^{\frac{1}{2}}Z^V_{\perp\perp},
\end{eqnarray}
where $Z^V_{J}$, $Z^V_{|\perp}$, and $Z^V_{\perp\perp}$ represent the corresponding renormalization factors arising from vertex corrections. Implicit in this notation is the dependence of the wave-function renormalization constant of the  ``heavy” auxiliary field from the color representations. 

The Wilson-line operators that define the spin-dependent potentials in terms of chromomagnetic field insertions into a Wilson loop \cite{Brambilla:2000gk,Pineda:2000sz} are also of interest in this study. Within the auxiliary-field formalism, each insertion of a chromomagnetic field into a Wilson loop can be related to the following current operator as, 
\begin{eqnarray}
\label{eq:currentoperatorspin}
J^{\mu\nu}_{F}(x) &=&\bar{h}_v(x)gF^{\mu\nu}_{\perp\perp}(x)h_v(x),
\end{eqnarray}
where $v=(1, 0,0,0)$ and the subscript $F$ indicates that the ``heavy” auxiliary fields are in the fundamental representation of the SU($3$)  group. The renormalization formula for this operator is , 
\begin{eqnarray}
J^{\mu\nu,\, \rm B}_{F}(x)&=&Z_{F}J^{\mu\nu,\, \rm R}_{F}(x) = Z_{A}^{1/2} Z_{h_v}Z^V_{F}J^{\mu\nu,\, \rm R}_{F}(x),
\end{eqnarray}
where $Z^V_{F}$ represents the renormalization factor stemming from the vertex corrections of the current $\bar{h}_v(x)gF^{\mu\nu}_{\perp\perp}(x)h_v(x)$.

%%%%%%%%%%%%%%%%%%%%%%%%%%%%%%%%%%%%%%%%%%%%%%%%
